\section{Empirical Design and Data}

The two return comparisons use different units of observation. To examine
the original spatial CAPM specification, we follow a fixed cohort and
summarize each issuer's subsequent returns by its median. To compare peer
rules after factor adjustment, we form a new roster each December and
observe one return per issuer per quarter in the following year.

The December 2015 formation ranks US common securities in CRSP CIZ by
market capitalization. Each selected share class has a security
identifier (PERMNO). We combine the classes belonging to the same
company at formation (PERMCO) into an issuer portfolio and follow
those classes over 2016--2025. The top 1,000
securities represent \ffnum[0]{quarterly-primary-population-issuers}
issuer portfolios before linkage and feature requirements. The
analytical formation roster contains
\ffnum[0]{quarterly-primary-roster} issuers, and
\ffnum[0]{quarterly-primary-issuers} enter the cross-sectional fit.
An issuer needs at least eight observed quarters and a converged
first-stage fit. A selected security remains eligible when it leaves
a later top-1,000 ranking.

For the fixed cohort, the outcome is the median of each issuer's observed
quarterly excess returns. Following \citet{fernandez2011}, we estimate its
market exposure by robust regression. Incomplete quarters and the treatment
of delistings determine which returns enter the median
(Appendix~\ref{app:data}).

The financial comparison begins with the annual histories of four ratios:
market equity over assets, market-to-book, dividend yield, and debt
maturity. A second input set supplies the dated financial quantities used
to construct those ratios, allowing the encoder to learn relationships
between their components. A third supplies a broader panel of eligible
financial fields from SEC CompanyFacts. Each January we represent an
issuer using the financial fields available
at formation. The encoder's context consists of historical examples labelled
with next-quarter market equity over assets, whose values are already
observed by the query date. Wide inputs are encoded in blocks and their
representations concatenated. Comparing each encoder with peers
formed directly from its own input set measures what encoding adds at
that information level. Filing cutoffs govern availability at formation,
although revised historical filings and checkpoints released after the
sample make this a retrospective exercise.

Three peer-selection rules can use exactly the same ratio histories.
The joint Fern\'andez distance compares annual ranks within each history
and combines the four distances by their root mean square. The direct-input
control imputes missing features with historical support medians, includes
missingness indicators, and standardizes the features before computing
distances. The encoder graph instead uses distances between the resulting
issuer representations. Each graph assigns weights to peers, whose weighted
returns enter the spatial lag.
Applying the same neighbour rule to all three distances isolates the choice
of representation. Holding a distance fixed while varying the weighting
rule measures how concentrating weight on fewer peers changes the fit.

For the repeated-formation panel the outcome is one issuer's excess return
in one quarter. With $y_t$ the vector of those returns, the spatial lag
model is
\[
  y_t=\alpha_t\mathbf{1}+X_t\gamma+\rho W_t y_t+u_t .
\]
Each column of $X_t$ multiplies an issuer's estimated factor exposure by
the corresponding quarterly factor return. The factor comparison considers
individual controls and joint specifications including the five
Fama--French factors (FF5) plus momentum. For the nonmarket factors we use
the study's compounded long--short proxies. Quarter effects absorb the common
level of returns. The matched sample requires observed outcomes and
converged exposures across the declared specifications, so a change in
fit cannot arise simply because a different set of issuers enters.
The estimated exposures use each issuer's full observed quarterly panel.
The factor comparison uses the joint Fern\'andez distance for its
characteristic graph. The encoder ranking instead compares each frozen
representation with its same-input control. Paired blocks of quarterly
likelihood contributions quantify variation conditional on the estimated
loadings and point fits.
